% Generated draft: requires mathematical and source review.
\documentclass[11pt]{article}
\usepackage[margin=25mm]{geometry}
\usepackage{fontspec}
\setmainfont{Latin Modern Roman}
\usepackage{amsmath,amssymb,mathtools,unicode-math}
\setmathfont{Latin Modern Math}
\usepackage{enumitem,needspace,xurl,hyperref}
\setlength{\parindent}{0pt}
\setlength{\parskip}{5pt}
\setlength{\emergencystretch}{3em}
\newcommand{\N}{\mathbb{N}}
\newcommand{\Z}{\mathbb{Z}}
\newcommand{\Q}{\mathbb{Q}}
\newcommand{\R}{\mathbb{R}}
\newcommand{\C}{\mathbb{C}}
\newcommand{\F}{\mathbb{F}}
\newcommand{\A}{\mathbb{A}}
\newcommand{\PP}{\mathbb{P}}
\newcommand{\E}{\mathbb{E}}
\newcommand{\eps}{\varepsilon}
\newcommand{\dd}{\,\mathrm{d}}
\begin{document}
\section*{UnsolvedMath ID: 20003183}
\textbf{Source problem code:} catalogue entry
\subsection{Definitions and mathematical statement}
Let $S$ be a closed connected oriented surface of genus $g\geq2$, and let $\rho:\pi_1(S)\to SL(n,\R)$ lie in a Hitchin component, namely a component containing principal Fuchsian representations. For each marked complex structure $J\in\mathcal T(S)$ the source invokes a unique $\rho$-equivariant harmonic map
\[
u_J:(\widetilde S,J)\longrightarrow X_n=SL(n,\R)/SO(n).
\]
The question is whether exactly one Teichm\"uller point $J$ makes it conformal. In the analytic formulation used by the aid, conformality means vanishing Hopf differential $\operatorname{Hopf}(u_J)=0$; any stronger unbranched-immersion requirement needs the relevant theorem. Structures are considered modulo marking-preserving isotopy, not as distinct tensors obtained by reparametrizing the same marked structure.

The aid attributes existence for every Hitchin representation to Labourie's proper harmonic-map energy theorem: the energy on $\mathcal T(S)$ attains a minimum whose Hopf differential vanishes. Uniqueness is reported for $n=2,3$, matching Wentworth's original comment citing [Lab06]. In contrast, the aid attributes to Sagman--Smillie the assertion that for every $g\geq3$ and $n\geq4$ there is a $PSL(n,\R)$ Hitchin representation with at least two inequivalent equivariant minimal maps. In the cited setting fixed-domain harmonic uniqueness forces these maps to determine distinct Teichm\"uller points. These are literature assertions whose complete references and precise equivalence conventions are not included.

The source uses $SL(n,\R)$, so lifting must be explicit. The kernel of $SL(n,\R)\to PSL(n,\R)$ is trivial for odd $n$ and $\{\pm I\}$ for even $n$, in either case contained in $SO(n)$. For even $n$, the central lifting obstruction lies in $H^2(S;\Z/2\Z)$ and is locally constant on representation components. A principal Fuchsian representation lifts, using a lift of the Fuchsian holonomy to $SL(2,\R)$ and the principal representation. Thus the obstruction vanishes throughout the connected Hitchin component. The symmetric space and its isometric action are unchanged by the central lift. Hence every accepted Sagman--Smillie example transfers to an $SL(n,\R)$ Hitchin counterexample with the same equivariant harmonic maps. The literal all-rank uniqueness assertion is thereby reported false, without treating the projective group as a silent replacement of the source group.

The aid reports genus-two high-rank uniqueness as open in the cited work, not independently certified current status. It also gives the restricted positive family $\rho=\iota_n\circ\rho_0$ of principal Fuchsian representations in every rank: the totally geodesic hyperbolic plane reduces the conformality condition to the unique marked uniformizing structure for $\rho_0$. Thus higher-rank counterexamples do not assert nonuniqueness for every Hitchin representation.

Original Comment 2.6 by Anna Wienhard asks the analogous question for maximal representations. No target Hermitian group, normalization, or theorem is supplied for that separate question, and the $SL(n,\R)$ counterexamples do not automatically settle it.
\subsection{Short English statement}
The source asks for uniqueness of the marked complex structure making a Hitchin harmonic map conformal. The supplied literature reports existence in general, uniqueness in ranks two and three, and higher-rank counterexamples for genus at least three. An explicit central-lifting argument preserves those examples in the source's exact $SL(n,\R)$ group. Genus two and the separate maximal-representation question retain their stated qualifications.
\subsection{Sources}
\begin{itemize}
\item Local registry: permanent catalogue entry, exact source code catalogue entry; \texttt{aim-topology-notes.json}, \texttt{notes[270]}; canonical tag ``question''.
\item AIM workshop \emph{Representations of Surface Groups} (2007), Question 2.4 attributed to Richard Wentworth and Francois Labourie, \url{https://aimath.org/WWN/surfacegroups/surfacegroups.pdf}.
\item Source reference [Cor88], Corollary 3.5, for fixed-domain harmonic uniqueness; Comment 2.5 by Richard Wentworth cites [Lab06], Theorem 9.3.2, for $n=2,3$. No bibliography entries or complete titles are supplied. Comment 2.6 by Anna Wienhard proposes maximal representations. The appended ``3 Surface Bundles'' is the next section heading.
\item Supplied research attempt 1, registry updated 14 August 2026: counterexample, refuted-at-run label, low novelty confidence, machine-generated and not peer reviewed. It attributes proper energy existence to Labourie and higher-rank examples to Sagman--Smillie, without full citations, and supplies the lifting and principal-Fuchsian reductions.
\item Verification identifies the official 2007 PDF but does not include further checked theorem statements.
\end{itemize}
\end{document}
